\section{Clustering}
In Clustering, we seek to learn groups or clusters from our input data, and in
so doing assign each instance to a cluster. Intuitively, we group objects such
that similar objects end up in the same group.
\subsubsection{Measures of Similarity}
\begin{figure}[h!]
  \centering
  \raisebox{-0.5\height}{\includegraphics[width=0.4\textwidth]{transitivity}}
  \hspace*{.2in}
  \raisebox{-0.5\height}{\includegraphics[width=0.4\textwidth]{proximity}}
\end{figure}
As measures of similarity, we can consider transitivity, where points are
clustered according to whether they are closely related to neighbours
regardless of how far apart the first and last neighbours are.\\
Alternatively, we may consider proximity, where points are clustered according
to some mutual shared distance, of which they cluster around a shared point.
\subsubsection{Evaluation}
Without labels of our data, we often are unable to evaluate how well our
predictions compare with the ground truth, and we are unsure of the number of
``clusters''.\\
It is hence not possible to have a clustering function that fits all kinds of
data, but for different problem settings we may be able to employ different
clustering algorithms:
\begin{itemize}
  \item \textbf{Partitional Algorithms}: Algorithms such as $k$-means clustering,
    Spectral Clustering techniques that divide data into a fixed number of
    clusters directly
  \item \textbf{Hierachical Algorithms}: Algorithms that build a hierarchy of clusters
    by either merging or splitting, such as Agglomerative Clustering techniques
\end{itemize}
These clustering algorithms allow us to focus on specific tasks at hand, such as data
exploration and data preprocessing.

\subsection{$k$-means Algorithm}
In $k$-means algorithm, instances that lie inside a cluster have small
inter-point distances (similarity by proximity). Assuming we are interested in
finding $k$ clusters, of which the $i$-th cluster is represented by a centroid
or prototype denoted $\mu^{[j]}\in \mathbb{R}^{m}$. We seek to find centroids
such that the sum of squares of distances to each instance to its closest
centroid is minimised.
\subsubsection{Problem Setting}
We have input data consisting of $n$-instances $ {\left\{ \mathbf{x}^{(i)}
\right\}}^{n}_{i=1}$ where $\mathbf{x}^{(i)}\in \mathbb{R}^{m}$.
For each instance, we create a set of indicator variables ${\left\{ \rho^{i[j]}
\right\}}^{k}_{j=1}$ where $\rho^{i[j]}\in \left\{ 0,1 \right\}$ for whether
$\mathbf{x}^{(i)}$ is assigned to cluster $j$. For each point
$\mathbf{x}^{(i)}$, let:
\[j^{*}=\argmin_l\left\lVert \mathbf{x}^{(i)}-\mathbf{\mu}^{[l]} \right\rVert_2^2.\] 
\[\rho^{i[j]}=\begin{cases}
  1,&\text{if }j=j^{*}\\
  0&\text{otherwise}
\end{cases}.\] 
Our objective function is defined as the sum of squared distances between each
instance and its assigned cluster:
\[\Theta=\sum_{i=1}^{n} \sum_{j=1}^{k} \rho^{i[j]}\left\lVert \mathbf{x}^{(i)}-\mathbf{\mu}^{[j]} \right\rVert_2^2.\]
This is known as the distortion function.
\subsubsection{Optimisation}
This function is hard to minimise: with ideal prototypes, we can solve by
assigning data points to the nearest cluster; with ideal assignments, we could
solve for the prototypes.
The $k$-means algorithm, or \textbf{Lloyd's Algorithm} is a greedy algorithm
that attempts to solve the optimisation problem:
\begin{itemize}
  \item Start by choosing a initial (random) value for each $\mu^{[j]}$
  \item \textbf{E-Step}: Minimise $\Theta$ w.r.t. ${\left\{ \rho^{i[j]}
    \right\}}^{n,k}_{i=1,j=1}$, keeping all $\mu^{[j]}$ constant.
    \begin{itemize}
      \item Each of the $n$ terms are independent and we can optimise each
        independently.
      \item We choose $p^{i[j]}=1$ for the value of $j$ giving the minimum
        value of the squared Euclidean distances:
        \[\rho^{i[j]}=\begin{cases} 1,&\text{if }j=\argmin_l\left\lVert
        \mathbf{x}^{(i)}-\mathbf{\mu}^{[l]} \right\rVert_2^2\\
      0&\text{otherwise}\end{cases}\]
    \end{itemize}
  \item \textbf{M-Step}: Minimise $\Theta$ w.r.t. ${\left\{ \mu^{[j]}
    \right\}}^{k}_{j=1}$, keeping all $\rho^{[j]}$ constant.
    \begin{itemize}
      \item $\Theta$ is a convex quadratic function of $\mu^{[j]}$. Taking the
        derivative and setting it to zero, we obtain:
        \[2 \sum_{i=1}^{n} \rho^{i[j]}\left( \mathbf{x}^{(i)}-\mathbf{\mu}^{[j]} \right).\] 
        \[\mathbf{\mu}^{[j]}=\frac{\sum_{i=1}^{n} \rho^{i[j]}\mathbf{x}^{(i)}}{\sum_{i=1}^{n} \rho^{i[j]}}.\] 
    \end{itemize}
  \item The E-Step and M-Step are repeated until convergence.\\
    \begin{proof}
      By definition,
      \[\Theta\left( \rho^{i[j]}_{(t)},\mathbf{\mu}^{[j]}_{(t)} \right)=
      \sum_{i=1}^{n} \sum_{j=1}^{k} \rho^{i[j]}_{(t)}\left\lVert \mathbf{x}^{(i)}-\mathbf{\mu}_{(t)}^{[j]} \right\rVert_2^2.\] 
      On the M-step:
      \[\Theta\left( \rho^{i[j]}_{(t)},\mathbf{\mu}^{[j]}_{(t)} \right)\le 
      \sum_{i=1}^{n} \sum_{j=1}^{k} \rho^{i[j]}_{(t)}\left\lVert \mathbf{x}^{(i)}-\mathbf{\mu}_{(t-1)}^{[j]} \right\rVert_2^2.\] 
      On the E-step:
      \[\sum_{i=1}^{n} \sum_{j=1}^{k} \rho^{i[j]}_{(t)}\left\lVert \mathbf{x}^{(i)}-\mathbf{\mu}_{(t-1)}^{[j]} \right\rVert_2^2\le \sum_{i=1}^{n} \sum_{j=1}^{k} \rho^{i[j]}_{(t-1)}\left\lVert \mathbf{x}^{(i)}-\mathbf{\mu}_{(t-1)}^{[j]} \right\rVert_2^2=\Theta\left( \rho^{i[j]}_{(t-1)},\mathbf{\mu}^{[j]}_{(t-1)} \right).\] 
      \[\therefore \Theta\left( \rho^{i[j]}_{(t)},\mathbf{\mu}^{[j]}_{(t)} \right)\le\Theta\left( \rho^{i[j]}_{(t-1)},\mathbf{\mu}^{[j]}_{(t-1)} \right)\qed\] 
    \end{proof}
    Hence the algorithm converges, but not necessarily to the global optimum.
\end{itemize}
\subsubsection{Limitations}
\begin{itemize}
  \item $k$ has to be specified prior to performing the algorithm.
  \item The object is non-convex, and the converged optimum is not necessarily global.
  \item Assumes clusters are hyperspheres, and works poorly on non-convex shaped clusters.
  \item Works poorly on clusters of different sizes.
  \item Sensitive to data-scaling
  \item Sensitive to noise and other small perturbations of data, due to hard
    assignments to clusters.
\end{itemize}
\subsection{Gaussian Mixture Model}
\subsubsection{Latent Variable Models}
\begin{definition}[Latent Variables]
  Variables that cannot be observed directly, and can only be inferred through
  a model.
\end{definition}
\begin{definition}[Latent Variable Model]
  A statistical model that relates a set of observable variables to a set of latent variables.
\end{definition}
Through usage of latent variable models, we are able to express joint
distributions more efficiently at the expense of model assumptions. We assume
conditional independence of the attributes on the latent variable.
\subsubsection{Mixture Models}
Not all random variables can be modelled solely with standard probability
distribution functions.  With a mixture model, we have access to a richer set
of probability distribution functions (often multimodal).
\subsubsection{Problem Setting}
Consider a set of unlabelled points ${\left\{ \mathbf{x}^{(i)}\in \mathbb{R}^{m}
\right\}}^{n}_{i=1}$ and $k$ clusters. Each point has an unknown latent cluster
assignment associated with it, $z^{(i)}\in \left\{ 1,\ldots,k \right\}$, which
is the outcome of a multinomial random variable $\mathcal{Z}$:
\[z^{(i)}\sim\text{Multinomial}\left( \pi^{[1],\ldots,\pi^{[k]}} \right).\] 
\[p_{\mathcal{Z}}\left( z^{(i)}=j \right) =\pi^{[j]},\quad \sum_{i=1}^{k}
\pi^{[i]}=1,\quad\pi^{[j]}\ge 0,\quad\forall j.\] 
Additionally, each $\mathbf{x}^{(i)}$ is the outcome of a Gaussian random
variable, $\mathcal{X}$:
\[\mathbf{x}^{(i)}|\left( z^{(i)}=j \right) \sim \mathcal{N}\left(
\mathbf{x}^{(i)};\mathbf{\mu}^{[j]},\mathbf{\Sigma}^{[j]} \right).\]
where $\mathbf{\mu}^{[j]}\in \mathbb{R}^{m},\mathbf{\Sigma}^{[j]}\in
\mathbb{R}^{m\times m}$ are the mean and covariance associated with cluster
$j$.\\
The joint distribution is hence given by:
\[p_{\mathcal{X}, \mathcal{Z}}\left( \mathbf{x}^{(i)},z^{(i)} \right)
=p_{\mathcal{X}}\left( \mathbf{x}^{(i)}|z^{(i)} \right) p_{\mathcal{Z}}\left(
z^{(i)} \right).\] 
\begin{align*}
  p_{\mathcal{X}}\left( \mathbf{x}^{(i)} \right) =&
  \sum_{z=1}^{k} p_{\mathcal{X}}\left( \mathbf{x}^{(i)}|z^{(i)} \right)
  p_{\mathcal{Z}}\left( z^{(i)} \right)\\
  =&\sum_{j=1}^{k} \mathcal{N}\left(\mathbf{x}^{(i)};\mathbf{\mu}^{[j]},\mathbf{\Sigma}^{[j]}\right)\pi^{[j]}
\end{align*}
The unconditional probability of each $\mathbf{x}^{(i)}$ is modelled as a
Gaussian Mixture Model.
\subsubsection{Optimisation}
We derive the log-likelihood:
\[\mathcal{L}=\sum_{i=1}^{n} \log\left( \pi^{[j]}\sum_{j=1}^{k}
\mathcal{N}\left( \mathbf{x}^{(i)};\mathbf{\mu}^{[j]},\mathbf{\Sigma}^{[j]}
\right) \right)\]
subject to:
\[\sum_{j=1}^{k} \pi^{[j]}=1,\quad\mathbf{\Sigma}^{[j]}\succ 0\] 
which is a constrained, non-concave optimisation problem.
Differentiating $L$ with respect to each variable, we have:
\begin{align*}
  \nabla_{\mathbf{\mu}^{[j]}} L=&\sum_{i=1}^{n} \frac{\partial L}{\partial \mathcal{N}\left(\mathbf{x}^{(i)};\mathbf{\mu}^{[j]},\mathbf{\Sigma}^{[j]}\right)} \frac{\partial \mathcal{N}\left(\mathbf{x}^{(i)};\mathbf{\mu}^{[j]},\mathbf{\Sigma}^{[j]}\right)}{\partial \mathbf{\mu}^{[j]}}\\
  =&\sum_{i=1}^{n} \frac{\pi^{[j]}}{\sum_{j=1}^{k} \pi^{[j]}\mathcal{N}\left(\mathbf{x}^{(i)};\mathbf{\mu}^{[j]},\mathbf{\Sigma}^{[j]}\right)}\frac{\partial \mathcal{N}\left(\mathbf{x}^{(i)};\mathbf{\mu}^{[j]},\mathbf{\Sigma}^{[j]}\right)}{\partial \mathbf{\mu}^{[j]}} \\
  =&\sum_{i=1}^{n} \frac{\pi^{[j]}\mathcal{N}\left(\mathbf{x}^{(i)};\mathbf{\mu}^{[j]},\mathbf{\Sigma}^{[j]}\right)}{\sum_{j=1}^{k} \pi^{[j]}\mathcal{N}\left(\mathbf{x}^{(i)};\mathbf{\mu}^{[j]},\mathbf{\Sigma}^{[j]}\right)}\frac{\partial}{\partial \mathbf{\mu}^{[j]}} \log\mathcal{N}\left(\mathbf{x}^{(i)};\mathbf{\mu}^{[j]},\mathbf{\Sigma}^{[j]}\right)\\
  =&-\sum_{i=1}^{n} \gamma^{i[j]}\frac{1}{2}\frac{\partial}{\partial \mathbf{\mu}^{[j]}} {\left( \mathbf{x}^{(i)}-\mathbf{\mu}^{[j]} \right) }^{T}\mathbf{\Sigma}^{-1}\left( \mathbf{x}^{(i)}-\mathbf{\mu}^{[j]}\right) \times \text{const.}\\
  =&\sum_{i=1}^{n} \gamma^{i[j]}\mathbf{\Sigma}^{-1}\left( \mathbf{x}^{(i)}-\mathbf{\mu}^{[j]}\right) \times \text{const.}
\end{align*}
where $\gamma^{i[j]}=\frac{\pi^{[j]}\mathcal{N}\left(\mathbf{x}^{(i)};\mathbf{\mu}^{[j]},\mathbf{\Sigma}^{[j]}\right)}{\sum_{j=1}^{k} \pi^{[j]}\mathcal{N}\left(\mathbf{x}^{(i)};\mathbf{\mu}^{[j]},\mathbf{\Sigma}^{[j]}\right)}$
\[\frac{\partial L}{\partial \Sigma^{[j]}}=\sum_{i=1}^{n} \gamma^{i[j]}\frac{\partial}{\partial \mathbf{\Sigma}^{[j]}} \log\mathcal{N}\left(\mathbf{x}^{(i)};\mathbf{\mu}^{[j]},\mathbf{\Sigma}^{[j]}\right).\]
Using Lagrange Multipliers to enforce the $\sum_{j=1}^{k} \pi^{[k]}=1$
constraint, we have:
\[\mathcal{L}=\sum_{i=1}^{n} \log\left( \sum_{j=1}^{k} \pi^{[j]}\mathcal{N}\left(\mathbf{x}^{(i)};\mathbf{\mu}^{[j]},\mathbf{\Sigma}^{[j]}\right) \right) +\lambda\left( \sum_{k=1}^{k} \pi_j-1 \right).\] 
\[\frac{\partial \mathcal{L}}{\partial \pi^{[j]}} =\sum_{i=1}^{n} \frac{\gamma^{i[j]}}{\pi^{[j]}}+\lambda.\] 
Solving for critical points of $L$ with respect to each variable, we have:
\[\mathbf{\mu}^{[j]*}=\frac{\sum_{i=1}^{n} \gamma^{i[j]}\mathbf{x}^{(i)}}{\sum_{i=1}^{n} \gamma^{i[j]}}.\] 
\[\mathbf{\Sigma}^{[j]*}=\frac{\sum_{i=1}^{n} \gamma^{i[j]}{\left(\mathbf{x}^{(i)}-\mathbf{\mu}^{[j]}\right)} {\left(\mathbf{x}^{(i)}-\mathbf{\mu}^{[j]}\right)}^{T}}{\sum_{i=1}^{n} \gamma^{i[j]}}.\] 
\[\pi^{[j]*}=\frac{1}{n}\sum_{i=1}^{n} \gamma^{i[j]}.\] 
\subsubsection{Responsibility}
$\gamma^{i[j]}$ can be viewed as the responsibility that component $j$ takes
for explaining $\mathbf{x}^{(i)}$:
\[\gamma^{i[j]}=\frac{\pi^{[j]}\mathcal{N}\left(\mathbf{x}^{(i)};\mathbf{\mu}^{[j]},\mathbf{\Sigma}^{[j]}\right)}{\sum_{j=1}^{k}
\pi^{[j]}\mathcal{N}\left(\mathbf{x}^{(i)};\mathbf{\mu}^{[j]},\mathbf{\Sigma}^{[j]}\right)}\]
Each optimal variable expression is in terms of the responsibility, and the
responsibility itself is a function of the optimal variables. 
\subsubsection{Expectation Maximisation}
Hence, we adopt an analoguue to the $k$-means algorithm to solve this problem.
\begin{itemize}
  \item Start by choosing a initial (random) value for
    ${\left\{\mathbf{\mu}^{[j]*},\mathbf{\Sigma}^{[j]*},\pi^{[j]*}\right\}}_{j=1}^{k}$
  \item \textbf{E-Step}: Compute the responsibilities $\gamma^{i[j]}$ given the
    current parameters
%    ${\left\{\mathbf{\mu}^{[j]*},\mathbf{\Sigma}^{[j]*},\pi^{[j]*}\right\}}_{j=1}^{k}$
%    \begin{itemize}
%      \item Each of the $n$ terms are independent and we can optimise each
%        independently.
%      \item We choose $p^{i[j]}=1$ for the value of $j$ giving the minimum
%        value of the squared Euclidean distances:
%        \[\rho^{i[j]}=\begin{cases} 1,&\text{if }j=\argmin_l\left\lVert
%        \mathbf{x}^{(i)}-\mathbf{\mu}^{[l]} \right\rVert_2^2\\
%      0&\text{otherwise}\end{cases}\]
%    \end{itemize}
  \item \textbf{M-Step}: Minimise $\mathcal{L}$ w.r.t.
    ${\left\{\mathbf{\mu}^{[j]*},\mathbf{\Sigma}^{[j]*},\pi^{[j]*}\right\}}_{j=1}^{k}$,
    keeping all $\gamma^{i[j]}$ constant.
\[\mathbf{\mu}^{[j]*}=\frac{\sum_{i=1}^{n} \gamma^{i[j]}\mathbf{x}^{(i)}}{\sum_{i=1}^{n} \gamma^{i[j]}}\qquad
\mathbf{\Sigma}^{[j]*}=\frac{\sum_{i=1}^{n} \gamma^{i[j]}{\left(\mathbf{x}^{(i)}-\mathbf{\mu}^{[j]}\right)} {\left(\mathbf{x}^{(i)}-\mathbf{\mu}^{[j]}\right)}^{T}}{\sum_{i=1}^{n} \gamma^{i[j]}}\qquad
\pi^{[j]*}=\frac{1}{n}\sum_{i=1}^{n} \gamma^{i[j]}.\] 
  \item The E-Step and M-Step are repeated until convergence.\\
    \begin{proof}
      At the end of the $t$-th epoch, we have (By Jensen's Inequality):
      \[\mathcal{L}\left( \theta_{(t)} \right) \le \sum_{i} \sum_{z^{(i)}} q_{\mathcal{Z}(t)}^{(i)}\left( z^{(i)} \right) \log\left( \frac{p_{\mathcal{X},\mathcal{Z}}\left( \mathbf{x}^{(i)},z^{(i)};\theta_{(t)} \right)}{q^{(i)}_{\mathcal{Z}(t)}\left( z^{(i)} \right) } \right) \]
      In the preceding M-step,
      \[\theta_{(t)}=\argmax_\theta\left(  q_{\mathcal{Z}(t)}^{(i)}\left( z^{(i)} \right) \log\left( \frac{p_{\mathcal{X},\mathcal{Z}}\left( \mathbf{x}^{(i)},z^{(i)};\theta_{(t)} \right)}{q^{(i)}_{\mathcal{Z}(t)}\left( z^{(i)} \right) } \right)\right).\] 
      \[\sum_{i} \sum_{z^{(i)}}  q_{\mathcal{Z}(t)}^{(i)}\left( z^{(i)} \right) \log\left( \frac{p_{\mathcal{X},\mathcal{Z}}\left( \mathbf{x}^{(i)},z^{(i)};\theta_{(t)} \right)}{q^{(i)}_{\mathcal{Z}(t)}\left( z^{(i)} \right) } \right)\ge\sum_{i} \sum_{z^{(i)}}  q_{\mathcal{Z}(t)}^{(i)}\left( z^{(i)} \right) \log\left( \frac{p_{\mathcal{X},\mathcal{Z}}\left( \mathbf{x}^{(i)},z^{(i)};\theta_{(t-1)} \right)}{q^{(i)}_{\mathcal{Z}(t)}\left( z^{(i)} \right) } \right)\] 
      RHS is the output of the preceding E-step, in which the tightest bound is achieved from
      setting $q_{\mathcal{Z}(t)}$:
      \[\sum_{i} \sum_{z^{(i)}}  q_{\mathcal{Z}(t)}^{(i)}\left( z^{(i)} \right) \log\left( \frac{p_{\mathcal{X},\mathcal{Z}}\left( \mathbf{x}^{(i)},z^{(i)};\theta_{(t-1)} \right)}{q^{(i)}_{\mathcal{Z}(t)}\left( z^{(i)} \right) } \right).\] 
    \[\mathcal{L}\left( \theta_{(t)}\right)\ge \mathcal{L}\left( \theta_{(t-1)} \right)\qed\]
    \end{proof}
    Similarly, the algorithm converges but not necessarily to the global optimum.
\end{itemize}
\subsubsection{Relationship with $k$-means}
Consider a GMM where we have `spherical' clusters, i.e.
$\mathbf{\Sigma}^{[j]}=\epsilon\mathbf{I},\forall j,\epsilon=$const.:
\[p_\mathcal{X}\left( \mathbf{x}^{(i)}|z^{(i)}=j \right) =\frac{1}{{\left( 2\pi\epsilon \right)}^{\frac{m}{2}} }\exp\left( -\frac{1}{2\epsilon}\left\lVert \mathbf{x}^{(i)}-\mathbf{\mu}^{[j]} \right\rVert_2^2 \right).\]
The responsibilities are given by:
\[\gamma^{i[j]}=\frac{\pi^{[j]}\exp\left( \frac{1}{2\epsilon}\left\lVert \mathbf{x}^{(i)}-\mu^{[j]}\right\rVert_2^2\right)}{\sum_{j} \pi^{[j]}\exp\left( \frac{1}{2\epsilon}\left\lVert \mathbf{x}^{(i)}-\mu^{[j]} \right\rVert_2^2\right)}\]
As $\epsilon\to 0$:
\begin{align*}
  \gamma^{i[j]}\to&\begin{cases}1,&\text{if }j=\argmin_g\left\lVert
  \mathbf{x}^{(i)}-\mathbf{\mu}^{[g]}
\right\rVert_2^2\\0&\text{otherwise}\end{cases}\\
      =&\rho^{i[j]}
\end{align*}
In this case, $pi^{[j]}$ no longer plays an active role and the only important
M-step update is $\mu^{[j]}$.\ i.e.\ as $\epsilon\to 0$, GMM clustering $\to$
$k$-means.
\subsubsection{Limitations}
With GMM, we are able to accommodate for less restrictive cluster shapes and
sizes than $k$-means (ellipsoids rather than circles). $k$ can be selected with
reference to validation set likelihood and is more robust than $k$-means.\\
The objective is non-convex and works poorly on non-convex clusters.

